\section{Bootstrap 近似估计量的抽样波动}
\label{sec:11-Mathematical-Statistics/06-Resampling-and-Model-Selection/01}

评审会上要交一个数：新排序模型把会话的中位停留时长从 \(83\) 秒抬到了 \(87\) 秒。有人问这 \(4\) 秒的误差范围是多少。

均值有 \(s/\sqrt n\)，比例有 \(\sqrt{\hat p(1-\hat p)/n}\)，回归系数有一整套现成公式。中位数之差没有。再往下走，AUC、固定精确率处的召回、两个模型指标的比值、一条带分桶和截断的自定义指标，它们的标准误要么推不出来，要么推出来的式子里带着未知的密度取值，代进去又叠一层估计。而报告里那个数如果不带波动范围，读它的人就只能凭感觉决定要不要上线。

上一单元把抽样分布讲成了一个思想实验，见\hyperref[sec:11-Mathematical-Statistics/01-Statistical-Models/03]{``抽样分布连接数据与推断''}：重复这次实验会得到的那批数构成一条曲线，而你手上只有曲线上的一个点。Bootstrap 提供的是把这个思想实验在数据上重放一遍的办法。

\subsection{它只替换了一样东西}

\begin{center}
\begin{tikzpicture}[x=1cm,y=1cm,font=\small,>=stealth]
  % ---- 左：真实世界 ----
  \node[font=\footnotesize,black!75] at (2.1,4.75) {真实世界：\(P\) 看不见};
  \draw[black!60] (1.3,3.75) rectangle (2.9,4.3);
  \node at (2.1,4.02) {\(P\)};
  \draw[->,red!70!black] (2.1,3.75) -- (2.1,3.15);
  \draw[->,dashed,black!40] (2.1,3.75) -- (0.7,3.15);
  \draw[->,dashed,black!40] (2.1,3.75) -- (3.5,3.15);
  \fill[red!70!black] (2.1,2.95) circle (2.6pt);
  \draw[black!45] (0.7,2.95) circle (2.6pt);
  \draw[black!45] (3.5,2.95) circle (2.6pt);
  \node[font=\footnotesize,black!60] at (2.1,2.5) {一次样本，加上没发生的那些};
  \draw[->,black!45] (2.1,2.25) -- (2.1,1.75);
  \draw[dashed,black!55,thick] plot[domain=-1.75:1.75,samples=60] ({2.1+\x},{0.85+0.85*exp(-\x*\x/0.75)});
  \draw[black!45] (0.2,0.85) -- (4.0,0.85);
  \fill[red!70!black] (2.1,0.85) circle (2.4pt);
  \node[font=\footnotesize,align=center,black!70] at (2.1,0.25) {\(\hat\theta\) 的抽样分布\\（量不出来）};
  % ---- 中间：替换 ----
  \draw[->,black!75,very thick] (4.15,4.02) -- (6.35,4.02);
  \node[font=\footnotesize,black!80,anchor=south] at (5.25,4.12) {把 \(P\) 换成 \(\hat P_n\)};
  % ---- 右：bootstrap 世界 ----
  \begin{scope}[shift={(5.6,0)}]
    \node[font=\footnotesize,black!75] at (2.1,4.75) {Bootstrap 世界：\(\hat P_n\) 在硬盘上};
    \draw[black!45] (1.15,3.75) -- (3.05,3.75);
    \foreach \sx in {1.25,1.55,1.9,2.35,2.6,2.95}
      \draw[blue!55!black,thick] (\sx,3.75) -- (\sx,4.22);
    \node[font=\footnotesize,anchor=west] at (3.1,4.0) {\(\hat P_n\)};
    \draw[->,blue!60!black] (2.1,3.7) -- (2.1,3.15);
    \draw[->,blue!60!black] (2.1,3.7) -- (0.7,3.15);
    \draw[->,blue!60!black] (2.1,3.7) -- (3.5,3.15);
    \fill[blue!60!black] (2.1,2.95) circle (2.6pt);
    \fill[blue!60!black] (0.7,2.95) circle (2.6pt);
    \fill[blue!60!black] (3.5,2.95) circle (2.6pt);
    \node[font=\footnotesize,black!60] at (2.1,2.5) {\(B\) 份重抽样，全都能真的算};
    \draw[->,black!45] (2.1,2.25) -- (2.1,1.75);
    \draw[blue!60!black,very thick] plot[domain=-1.75:1.75,samples=60] ({2.1+\x},{0.85+0.85*exp(-\x*\x/0.75)});
    \draw[black!45] (0.2,0.85) -- (4.0,0.85);
    \node[font=\footnotesize,align=center,black!70] at (2.1,0.25) {\(\hat\theta^\star\) 的分布\\（跑一遍就有）};
  \end{scope}
\end{tikzpicture}

\emph{左右两侧的箭头结构一模一样，差别只在最上面那一格：一个抽不了，一个抽得了。}
\end{center}

图里两条流程从上到下逐格对应：一个分布，从它独立抽 \(n\) 个观测，把观测代进同一个 \(T\)，看结果散成什么样。方法的全部内容就是最上面那一次替换，用经验分布

\[
\hat P_n=\frac1n\sum_{i=1}^n\delta_{x_i}
\]

顶替未知的 \(P\)。\(\delta_x\) 是集中在点 \(x\) 上的单位质量，所以 \(\hat P_n\) 把总概率均分给 \(n\) 条观测，每条 \(1/n\)。从这样一个分布里独立抽 \(n\) 个点，具体做起来就是从原数据里有放回地抽 \(n\) 次。记

\[
x_1^\star,\ldots,x_n^\star\overset{\iid}{\sim}\hat P_n,
\qquad
\hat\theta^\star=T(x_1^\star,\ldots,x_n^\star),
\]

重复 \(B\) 次得到 \(\hat\theta^{\star1},\ldots,\hat\theta^{\star B}\)，它们在给定原数据下的条件分布，就是图右下角那条曲线，用来近似左下角那条量不出来的曲线。

把这次替换记牢，后面判断 bootstrap 何时失手就有了统一的入口：凡是 \(\hat P_n\) 离 \(P\) 还远、或者 \(T\) 对这点差别格外敏感的场合，右边那条曲线就不再像左边。

\subsection{有放回与样本量不变都不是可调项}

这两条常被当作实现细节，其实各自被上面那次替换钉死了。

有放回是从 \(\hat P_n\) 独立抽样的直接后果。每条观测的概率都是 \(1/n\)，独立抽 \(n\) 次自然允许同一条被抽到两遍、另一条一次都不出现。若改成无放回地抽 \(n\) 个，每次得到的只是原数据的一个排列；对均值、中位数这类与顺序无关的统计量，\(B\) 次重抽会给出 \(B\) 个完全相同的数，曲线塌成一根竖线。想看的是数据换一批估计会怎么变，而排列并没有换数据。

保持样本量为 \(n\)，是为了让被模仿的问题保持同一个尺度。真实问题是``从 \(P\) 抽 \(n\) 个观测算 \(\hat\theta\)''，重抽样模仿的是同一句话，只把 \(P\) 换掉。抽 \(n/2\) 个点，量到的是样本量减半时的波动，比真实标准误大出约 \(\sqrt2\) 倍。有一类 \(m\)-out-of-\(n\) 变体确实取 \(m<n\)，那是为了修复参数落在边界这类非正则情形而有意为之，与随手改小样本量是两回事。

\subsection{标准误和偏差都是这条曲线的读数}

有了 \(B\) 个重抽样统计量，标准误就是它们的样本标准差，

\[
\widehat{\SE}_{\mathrm{boot}}
=\sqrt{\frac1{B-1}\sum_{b=1}^B\bigl(\hat\theta^{\star b}-\bar\theta^\star\bigr)^2},
\qquad
\bar\theta^\star=\frac1B\sum_{b=1}^B\hat\theta^{\star b},
\]

偏差则由重抽样均值与原估计之差给出，

\[
\widehat{\Bias}_{\mathrm{boot}}=\bar\theta^\star-\hat\theta.
\]

两个量都是条件于已观测数据算出来的，它们描述的是图右下角那条曲线的宽度和中心偏移。

拿样本均值对一遍账。条件于数据，\(x_i^\star\) 从 \(\hat P_n\) 独立抽取，其方差就是经验分布的方差，于是

\[
\begin{aligned}
\Var^\star(\bar x^\star)
&=\frac1n\Var^\star(x_1^\star)\\
&=\frac1n\cdot\frac1n\sum_{i=1}^n(x_i-\bar x)^2\\
&=\frac{n-1}{n}\cdot\frac{s^2}{n}.
\end{aligned}
\]

结果是经典答案 \(s^2/n\) 乘上 \((n-1)/n\)，差的那一点来自经验分布用 \(1/n\) 而非 \(1/(n-1)\) 做归一。对账一下用到了什么：第一步只用了条件独立与样本量为 \(n\)，第二步只用了 \(\hat P_n\) 的定义，全程没有碰 \(P\) 的形状。所以这个等式在任何数据上都成立，它验证的是重抽样机制本身没错，不构成 bootstrap 对复杂统计量也有效的证据。均值这里 bootstrap 没带来新东西；它的价值在于把同样几步原样搬到中位数之差、AUC 或一整条打分流水线上，而那些地方的解析推导已经不可行。

置信区间还需要多走一步，且几种做法并不等价。正态区间取 \(\hat\theta\pm z_{1-\alpha/2}\widehat{\SE}_{\mathrm{boot}}\)，简单对称，遇到偏斜的抽样分布覆盖率就不对。分位数区间直接取 \([q^\star_{\alpha/2},\,q^\star_{1-\alpha/2}]\)，让区间自动跟随曲线的偏斜，前提是它默认重抽样分布的形状已经忠实复制了抽样分布的形状。基本区间把 \(\hat\theta^\star-\hat\theta\) 的分布当作 \(\hat\theta-\theta\) 的镜像，因而写成 \([2\hat\theta-q^\star_{1-\alpha/2},\,2\hat\theta-q^\star_{\alpha/2}]\)；镜像是一个近似，不是恒等式，参数有边界时它可能给出越界的端点。BCa 再对偏差和加速度做修正，覆盖阶数更高，但实现复杂，对小样本也更敏感。用了 bootstrap 并不等于区间覆盖正确，报告时把区间类型写出来，读的人才知道你担保了什么。

\subsection{失效发生在经验分布够不着的地方}

先看一个能一眼看穿的例子。设 \(x_1,\ldots,x_n\overset{\iid}{\sim}\Uniform(0,\theta)\)，估计量取样本最大值 \(\hat\theta=\max_i x_i\)。真实抽样分布里，\(\hat\theta\) 有一侧的波动来自``下次可能抽到更大的值''。而任何重抽样样本都只能从已有的 \(n\) 条里挑，它的最大值不可能超过 \(x_{(n)}\)，右侧波动被整齐地切掉了。更糟的是，重抽样恰好抽中最大值那条的概率是

\[
P^\star\bigl(\hat\theta^\star=\hat\theta\bigr)=1-\Bigl(1-\frac1n\Bigr)^{n}\xrightarrow[n\to\infty]{}1-e^{-1}\approx0.632,
\]

所以曲线在右端立着一根占了六成质量的柱子，形状与真实抽样分布毫无关系。

这个例子的价值在于它的症状是可以观察的。把 \(B\) 个重抽样统计量画成直方图，如果右端（或左端）出现一根异常高的柱子，或者分位数区间的某个端点恰好等于 \(\hat\theta\) 本身、把 \(B\) 从 \(1000\) 加到 \(10000\) 也纹丝不动，那说明这一侧的波动根本没被生成出来，而不是它真的很小。

小样本失效的症状类似但更常见。\(n=20\) 时重抽样统计量只能落在由 \(20\) 个原始值组合出来的有限个点上，直方图是几根离散的柱子，估计 \(90\%\) 分位数时更是几乎只在两三条观测之间跳。此时区间端点往往正好等于某一条原始观测值，加大 \(B\) 只会让这个端点更稳定地等于它，而稳定不等于正确。判断办法是把重抽样统计量的不同取值个数数一遍：如果 \(1000\) 次重抽只产生了几十个互异的值，曲线的分辨率已经不足以支撑分位数。

重尾数据的症状则是少数点说了算。若数据来自尾指数很小的分布，样本方差本身就不稳定，\(\widehat{\SE}_{\mathrm{boot}}\) 会被最大的一两条观测主导：含有那条记录的重抽样样本给出一个数，不含的给出另一个数，直方图分裂成双峰。一个便宜的诊断是逐条留一重算 \(\hat\theta\)，看最大的那次变动有多大；若删掉单条记录就能让 \(\hat\theta\) 移动一个 \(\widehat{\SE}_{\mathrm{boot}}\) 的量级，说明估计由个别点决定，重抽样给出的宽度不可信。参数落在边界上的情形（例如方差分量的估计恰为零）会留下另一种痕迹：重抽样分布在边界处堆积出一个原子，标准误被机械地压小。

还有两类失效跟数据本身无关，跟你怎么用它有关。其一是重抽单位选错：时间序列、同一用户的多条行为、同一门店的多笔订单，逐行有放回重抽等于宣称这些记录相互独立，得到的标准误往往只有真实值的一半甚至更小。补救办法是按依赖单位重抽——时间序列切成长度为 \(\ell\) 的块做块 bootstrap，聚类数据以聚类为单位整块抽出。选择的重抽单位就是你对独立性的声明，它必须与数据的生成方式对得上。其二是把一部分建模步骤留在了循环外面：先用全量数据选了变量，再对选定的模型做 bootstrap，得到的曲线只反映系数估计的波动，完全看不见``换一批数据会选出另一组变量''这件事，区间因此系统性偏窄。凡是用到了数据的步骤都得搬进重抽样循环里重做一遍，这条原则在下下节讲交叉验证时会以更尖锐的形式再出现。

\subsection{\(B\) 只买 Monte Carlo 精度}

\(B\) 容易被当成万能旋钮：结果不好看，就多跑几轮。区分清楚三层误差就知道它买不到什么。最外面一层是原样本的抽样误差，你只有 \(n\) 条观测，这批观测本身可能不具代表性；中间一层是用 \(\hat P_n\) 顶替 \(P\) 带来的误差，它由 \(n\) 决定；最里面一层才是只跑了 \(B\) 次而非无穷次带来的 Monte Carlo 噪声。加大 \(B\) 只让第三层变小，前两层纹丝不动。换句话说，\(B=10^5\) 的 bootstrap 会非常精确地告诉你，一个可能有偏的近似给出了什么答案。

按经验，标准误估计在 \(B\) 取两三百时已经稳定，分位数区间通常需要 \(1000\) 以上，尾部分位数还要更多。判断是否够用有个直接办法：换一个随机种子重跑，看区间端点变动多少，如果变动已经小到不影响结论就够了。报告时把 \(B\)、随机种子、重抽单位和区间类型一并写上，别人才有可能复现这个区间。

Bootstrap 回答的是``这个数会晃多大''。下一节换一个问题：两组之间的差距是不是可以完全由分组的偶然性解释。那里的重排动的不是观测本身而是标签，依据也从``经验分布接近真实分布''换成了另一条更强、也更容易被现实数据破坏的假设。
